\chapter*{ZAKOŃCZENIE}
\addcontentsline{toc}{chapter}{\protect\hspace{0.45cm} ZAKOŃCZENIE}

Niniejsza praca magisterska stanowiła całościowe studium nad rolą i~wpływem wieloetapowych procedur czyszczenia oraz wstępnego przetwarzania danych tabelarycznych na jakość predykcji i~stabilność algorytmów uczenia maszynowego. Wszystkie sformułowane we wstępie cele badawcze oraz pytania problemowe zostały w~pełni zrealizowane i~rozstrzygnięte w~toku przeprowadzonych badań.

W~pracy wykazano, że trenowanie modeli na danych obarczonych zanieczyszczeniami (brakami danych, skrajnymi anomaliami oraz szumem kategorycznym) prowadzi do drastycznego załamania ich skuteczności decyzyjnej. Jednocześnie udowodniono, że odpowiednio skonstruowany potok przygotowania danych pozwala na niemal całkowite zniwelowanie negatywnych skutków degradacji danych wejściowych, zabezpieczając wysoką jakość klasyfikacji nawet w~warunkach skrajnego, 50-procentowego zanieczyszczenia próby.

Zgromadzony materiał badawczy oraz przeprowadzona weryfikacja statystyczna pozwoliły na jednoznaczne potwierdzenie kluczowej hipotezy pracy: to sekwencyjne połączenie filtracji wartości odstających, imputacji braków oraz standaryzacji cech przynosi efekt synergii, którego nie sposób osiągnąć stosując wyłącznie pojedyncze, izolowane techniki czyszczenia. Ponadto badania uwidoczniły istotną asymetrię wrażliwości pomiędzy poszczególnymi architekturami klasyfikatorów, dowodząc, że dobór procedur przygotowania danych musi uwzględniać matematyczną specyfikę i~założenia leżące u~podstaw danego modelu.

Podsumowując, uzyskane rezultaty stanowią spójną klamrę dla podjętej tematyki badawczej i~jednoznacznie wskazują, że jakość danych zasilających algorytmy sztucznej inteligencji stanowi fundamentalny warunek budowy wiarygodnych oraz stabilnych systemów predykcyjnych.
